\documentclass{article}

\usepackage{amsmath, amssymb}

% Macro: shorthand for the risk-neutral measure.
\newcommand{\Q}{\mathbb{Q}}

\title{My First LaTeX Document}
\author{Your Name}
\date{\today}

\begin{document}
\maketitle

\section{Introduction}
This document demonstrates the basic LaTeX workflow. In later sections of the course we will produce data and figures from a data pipeline, as well as adding references from a \verb|.bib| file.

\section{Math: Inline and Displayed}
Hello, world. The Pythagorean theorem says
\begin{equation}
  a^2 + b^2 = c^2.
\end{equation}
We can also write math inline: a right triangle with legs $a = 3$ and $b = 4$ has hypotenuse $c = 5$.

\section{A Sliver of Asset Pricing}
Under the risk-neutral measure $\Q$, the price of a payoff $X_T$ at date $T$ is
\begin{equation}
  P_0 = E^{\Q}\left[ e^{-\int_0^T r_s\,ds} \, X_T \right].
\end{equation}
The macro \verb|\Q| expands to $\Q$ --- every reference to the risk-neutral measure uses the same symbol, and if we ever change notation we change it in one place.
\end{document}
